\documentclass[11pt,a4paper]{article}
\usepackage[utf8]{inputenc}
\usepackage{amsmath}
\usepackage{amsfonts}
\usepackage{graphicx}
\usepackage[table,xcdraw]{xcolor}


\usepackage{caption}
\usepackage{subcaption}

\renewcommand{\familydefault}{\sfdefault} % cambiamos la fuente a una sans

\usepackage{float} % para que floten las imagenes o algo asi...
\usepackage{wallpaper} %paquete para usar una imagen como encabezado!
\usepackage{hyperref} %para usar hypervinculos 
\usepackage[export]{adjustbox} %para usar marcos en imagenes
\usepackage{eurosym} % para el euro
\usepackage{transparent} %para las marcas de agua
\usepackage{eso-pic}  %para las marcas de agua

\definecolor{azul_marcos}{RGB}{0,128,159} %defino el color azul de los marcos
\usepackage{sectsty} %esto es para cambiar el color de las fuentes creo
\renewcommand{\familydefault}{\sfdefault} % cambiamos la fuente a una sans
\sectionfont{\color{azul_marcos}}  % sets colour of sections
\subsectionfont{\color{azul_marcos}}  % sets colour of sections


\usepackage{pdfpages} %para insertar pdfs
\usepackage{amssymb}
\usepackage{pstcol} % para color
\usepackage{pst-node} % para diagramas
\usepackage{pst-plot} % para representacion de dat
\usepackage[spanish]{babel}
\addto\captionsspanish{\renewcommand\chaptername{Bloque}}
%\usepackage[total={18cm,21cm},top=2cm, left=2cm]{geometry}
\usepackage{anysize}
\pagestyle{plain}
%\markboth{left head}{right head}
%\markright{Guía de impresión FlexiSMART}
\marginsize{3cm}{2cm}{2.5cm}{1cm}
\title{Guía de impresión FlexiSMART}
\date{}

%configuracion de la marca de agua
\AddToShipoutPicture{
    \put(0,0){
        \parbox[b][\paperheight]{\paperwidth}{%
            \vfill
            \centering
            {\transparent{0.2}\includegraphics[scale=1.25]{FOTOS/logofff}}%
            \vfill
        }
    }
}

\begin{document}
\ULCornerWallPaper{1}{FOTOS/header}
\LLCornerWallPaper{1}{FOTOS/footer}
%\maketitle
%\tableofcontents
\includepdf{PDF/ES_PORTADA.pdf}
\section{<%=strings[1]%>}<%=strings[2]%> 
\section{<%=strings[3]%>}
<%=strings[4]%>
\\\\
<%=strings[5]%>
\\\\
<%=strings[6]%>
\\\\
<%=strings[7]%>
\\\\
<%=strings[8]%>
\section{Ficha técnica y parámetros de impresión}

\begin{table}[H]
\centering
\caption*{Ficha técnica}
\begin{tabular}{|
>{\columncolor[HTML]{FFFFFF}}l |
>{\columncolor[HTML]{FFFFFF}}c |}
\hline
\multicolumn{1}{|c|}{\cellcolor[HTML]{FFFFFF}\textbf{Material}}   & FlexiSMART (TPE)   \\ \hline
\textbf{Colores disponibles}              & 11                 \\ \hline
\textbf{Formatos disponibles}             & 1kg, 250gr         \\ \hline
\textbf{Temperatura de deflexión térmica} & 90ºC               \\ \hline
\textbf{Temperatura de fusión}            & 160ºC              \\ \hline
\textbf{Temperatura de descomposición}    & \textgreater 240ºC \\ \hline
\textbf{Densidad}                         & 0.96 gr / cm3      \\ \hline
\textbf{Estiramiento máximo}              & 600\%              \\ \hline
\end{tabular}
\end{table}
\begin{table}[H]
\centering
\caption*{Parámetros de impresión recomendados usando un nozzle de 0.4 mm}
\begin{tabular}{|
>{\columncolor[HTML]{FFFFFF}}l |
>{\columncolor[HTML]{FFFFFF}}c |}
\hline
\multicolumn{1}{|c|}{\cellcolor[HTML]{FFFFFF}\textbf{Temperatura recomendada de impresión}} & 195º-220º              \\ \hline
\textbf{Velocidad recomendada de impresión}                         & 20-60mm/s              \\ \hline
\textbf{Temperatura cama caliente}                                  & \textgreater 18º        \\ \hline
\textbf{Altura de capa óptima}                                      & 0.2 mm                 \\ \hline
\textbf{Perímetros}                                                 & 3                      \\ \hline
\textbf{Top solid layers}                                           & 5                      \\ \hline
\textbf{Retracción}                                                 & Desactivada o reducida \\ \hline
\end{tabular}
\end{table}

Puedes descargarte nuestros perfiles completos de impresión de los principales programas de laminación (Cura, Slic3r y Simplify3D) desde nuestra página web:
\\\\
\centerline{ {\huge \url{www.fffworld.com/documentation} } }
\\\\
Los parámetros óptimos dependerán de la impresora 3D que utilices, sin embargo, son unos buenos parámetros para tenerlos como punto de partida. Con unas pocas impresiones serás capaz de encontrar los límites y la configuración perfecta para tu maquina.
\section{<%=strings[9]%>}
	\subsection{<%=strings[10]%>}
<%=strings[11]%>
\\\\
<%=strings[12]%>
		\subsubsection{<%=strings[13]%>}
<%=strings[14]%>
\\\\
<%=strings[15]%>
\\\\
<%=strings[16]%>
\\\\
<%=strings[17]%>
		\subsubsection{<%=strings[18]%>}
<%=strings[19]%>
\\\\
<%=strings[20]%>
%FOTO BARREL
	\subsection{<%=strings[21]%>}
<%=strings[22]%>
\\\\
<%=strings[23]%>
	\begin{description}
		\item[<%=strings[24]%>] <%=strings[25]%> 
		\item[<%=strings[26]%>] <%=strings[27]%> 
	\end{description}
	\subsection{<%=strings[28]%>}
<%=strings[29]%>
\\\\
<%=strings[30]%>
	\subsection{<%=strings[31]%>}
<%=strings[32]%>
\\\\
<%=strings[33]%>
\\\\
<%=strings[34]%>
\\\\
<%=strings[35]%>
\\\\
<%=strings[36]%>
	\subsection{<%=strings[37]%>}
<%=strings[38]%>
\\\\
<%=strings[39]%>
\\\\
<%=strings[40]%>
\\\\
<%=strings[41]%>
	\subsubsection{<%=strings[42]%>}
<%=strings[43]%>
\\\\
<%=strings[44]%>
\\\\
<%=strings[45]%>
	\subsubsection{<%=strings[46]%>}
<%=strings[47]%>
\\\\
<%=strings[48]%>
\\\\
<%=strings[49]%>
\\\\
<%=strings[50]%>
\url{https://www.antonmansson.com/how-to-cold-pull-clogged-nozzle/}
\url{https://www.trideus.be/en/blogs/stories/tips-tricks-do-the-cold-pull/}
\url{https://ultimaker.com/en/resources/19510-how-to-apply-atomic-method}
\url{https://printrbot.zendesk.com/hc/en-us/articles/202100554-Unclogging-the-Hot-End-Using-the-Cold-Pull-Method}
	\subsubsection{<%=strings[51]%>}
<%=strings[52]%>
\\\\
<%=strings[53]%>
\\\\
<%=strings[54]%>
\\\\
<%=strings[55]%>
\section{¿Quieres apoyar nuestro proyecto?}
Todos los miembros FFF World amamos la impresión 3D y la comunidad maker. Nos sentimos afortunados de poder trabajar en proyectos donde podamos entregar nuestra pasión sincera. En el futuro, nos gustaría poder desarrollar más materiales, más colores, más formatos. En definitiva, nos gustaría poder hacer crecer nuestra empresa.
\\\\
Para ello, una de las principales acciones para ayudarnos, si quieres hacerlo y estás satisfecho con el filamento, es la de votarnos en Amazon con 5 estrellas.
\begin{figure}[H]
\centering
\includegraphics[width=0.5\textwidth,cfbox=azul_marcos 1pt 0pt]{FOTOS/AMAZON_FIVE_STARS}
\caption*{¡Muchas gracias!}
\end{figure}
\subsection{Another filaments with outstanding properties available now in Amazon}
\begin{description}
\item[FlexiSMART Tech:] Diseñado para resistir a la abrasión y al desgaste de impresiones técnicas.
\item[ABS Tech:] Efecto warping minimizado. Alto rendimiento en aplicaciones técnicas.
\item[PETG Tech:] Máxima resistencia mécanica. Resistente al contacto con el agua y los rayos UV. Apto para uso alimentario.
\item[FilaMETAL:] PLA con carga metálica no abrasiva que da un acabado metálico espectacular a tus impresiones.
\item[PC Tech:] Policarbonato con gran resistencia a la temperatura y con excelentes propiedades mecánicas.
\item[Nylon Tech:] Imprimible a baja temperatura. Resistencia a los golpes con cierto grado de flexibilidad.
\item[PVA Tech:] Filamento soluble en agua indicado para uso como material de soporte. Excelente compatibilidad con PLA.
\item[HIPS Tech:] Filamento soluble en limoneno indicado para uso como material de soporte. Buena resistencia mecánica y excelente compatibilidad con ABS.
\end{description}
\includepdf{PDF/ES_CONTRAPORTADA.pdf}
\end{document}
